% 五个预测残差符号的Huffman树，概率为教学构造(0.4,0.2,0.2,0.1,0.1)；并列时按正文约定合并。宽约90mm。
\begin{tikzpicture}[x=1cm,y=1cm,font=\figurefont\small,
  leaf/.style={draw=FigureBlue,fill=FigureBlueFill,rounded corners=1.8pt,
    minimum width=13mm,minimum height=7mm,align=center,inner sep=1pt},
  inner/.style={draw=FigureYellow,fill=FigureYellowFill,circle,
    minimum size=8mm,inner sep=0pt},
  edge/.style={draw=BookMuted,line width=.8pt},
  bit/.style={text=BookTeal,font=\figurefont\footnotesize}]
 \node[inner] (r) at (3.2,3.6) {$1$};
 \node[leaf] (z) at (1.2,2.4) {$0$\\$0.4$};
 \node[inner] (n3) at (5.2,2.4) {$0.6$};
 \node[inner] (n2) at (3.6,1.2) {$0.4$};
 \node[inner] (n1) at (6.8,1.2) {$0.2$};
 \node[leaf] (p1) at (2.8,0) {$+1$\\$0.2$};
 \node[leaf] (m1) at (4.4,0) {$-1$\\$0.2$};
 \node[leaf] (p2) at (6.0,0) {$+2$\\$0.1$};
 \node[leaf] (m2) at (7.6,0) {$-2$\\$0.1$};
 \draw[edge] (r)--node[bit,above left] {0} (z);
 \draw[edge] (r)--node[bit,above right] {1} (n3);
 \draw[edge] (n3)--node[bit,above left] {0} (n2);
 \draw[edge] (n3)--node[bit,above right] {1} (n1);
 \draw[edge] (n2)--node[bit,left] {0} (p1);
 \draw[edge] (n2)--node[bit,right] {1} (m1);
 \draw[edge] (n1)--node[bit,left] {0} (p2);
 \draw[edge] (n1)--node[bit,right] {1} (m2);
 \node[text=BookMuted,font=\figurefont\footnotesize,anchor=north] at (1.2,1.85) {码字 0};
 \node[text=BookMuted,font=\figurefont\footnotesize,anchor=north] at (5.2,-.55) {码字 100, 101, 110, 111};
\end{tikzpicture}
